\subsection{Circuitos variacionales y circuitos aleatorios: entre regularidad y alta entropía}
\label{subsec:patrones_variacionales_aleatorios}

En aplicaciones NISQ, algunos circuitos variacionales pueden heredar regularidades de la función objetivo. En esos casos cabe esperar distribuciones de amplitud menos uniformes que las de un circuito completamente aleatorio, por ejemplo con subconjuntos dominantes o con simetrías parciales, lo que en principio favorecería más la compresión que un estado de alta entropía.

En cambio, para circuitos aleatorios profundos lo más razonable es esperar el escenario opuesto: estados densos, con gran diversidad de amplitudes y poca repetición exacta. Bajo esa condición, una estrategia sin pérdida tendría poco margen para explotar redundancia.
